\documentclass{article}
\usepackage{amsmath}
\usepackage{amssymb}
\usepackage{array}
\usepackage{booktabs}
\usepackage{textcomp}
\usepackage{graphicx}

\usepackage{amsmath,amsfonts,amsthm,tikz} %packages for math symbols
\newtheorem{theorem}{Theorem}[section]
\newtheorem{lemma}{Lemma}[section]

\newenvironment{solution}
  {\begin{proof}[Solution]}
  {\end{proof}}


\title{hw 1}
\author{Ian Chi}
\date{\today}
\begin{document}
\maketitle

\begin{enumerate}
    \item \begin{enumerate}
        \item the sum over all x and y should add up to 1. 
        \begin{equation}
            \sum_{x\in\{0,1,2\}}^{}\sum_{y\in\{0,1\}}^{}c(x+y) = 1
        \end{equation}

        this implies \(c =\left( \sum_{x\in\{0,1,2\}}^{}\sum_{y\in\{0,1\}}^{}(x+y) \right)^{-1}\)
        \(c = 1/9\)


        \item \begin{equation}
            P(x) = \sum_{y}^{} c(x+y) = c(x+1)
        \end{equation}
        \begin{equation}
            P(y) = \sum_{x}^{} c(x+y) = c(y+3)
        \end{equation}
        \item \begin{equation}
            P(x|y=1) = \sum_{x}^{}c_2(x+1)
        \end{equation}
        where \(c_2 = 1/6\)

        \item no. the distribution for x changes based on whether y is 0 or 1.


    \end{enumerate}

    \item \begin{enumerate}
        \item \begin{equation}
            \begin{split}
                P_{z_i}(1) = p\\
                P_{z_i}(0) = 1-p\\
            \end{split}
        \end{equation}
        \item the pmf just follows a binomial distribution with \(p=p\)
        \begin{equation}
            P(x)  = {n\choose x} p^x (1-p)^{n-x}
        \end{equation}
        \begin{equation}
            P(x\geq k) = \sum_{k}^{n} {n\choose x} p^x (1-p)^{n-x}
        \end{equation}
        \item \begin{equation}
            P(Z_1 = 1|S_n\geq 1) =\frac{P(S_n\geq 1| Z_1 = 1) P(Z_1)}{P(S_n\geq 1)} = \frac{1 \cdot p}{\sum_{k}^{n} {n\choose x} p^x (1-p)^{n-x}}
        \end{equation}
        the probability changes since we know one at least is true, each one is more likely to be working.
        \item \begin{equation}
            P(T > m | S_n \geq 1) = \frac{1}{(1-p)^T}P(Z_1 = 1|S_{n-m}\geq 1)
        \end{equation}
        this question is the same as the original problem except the first m failed. the probability that the first m failed is \(\frac{1}{(1-p)^T}\) thus we multiply them together and they are independent since no packet affects the outcome of the next packet.
    \end{enumerate}
    \item \begin{enumerate}
        \item the distribution is an exponential distribution but discrete. 
        \begin{equation}
            P(T=k) = (1-p)^{k-1}p
        \end{equation}
        \item \begin{equation}
            1- \sum_{i=0}^k  (1-p)^{i-1}p = (1-p)^k
        \end{equation}
        \item apply \begin{equation}
            \sum kr^{k-1} = \frac{1}{(1-r)^2}
        \end{equation} for r between 0 and 1. since 1-p is between 0 and 1,\begin{equation}
            \mathbb E(T) = \sum kP(T=k) = \frac{1}{p}
        \end{equation}
        \item 
        \begin{equation}
            \begin{split}
                P(T>m+n|T>m) &= \frac{P(T>m|T>m+n)P(T>m+n)}{P(T>m)} \\
                &= \frac{1\cdot(1-p)^{m+n}}{(1-p)^m} = (1-p)^n
            \end{split}
        \end{equation}

        it means we sort of can restart our algorithm at the beginning if something fails before. 

    \end{enumerate}
    \item \begin{enumerate}
        \item the probability is \(\frac{1}{100}\). 
        \item the probability is concentrated into one door and thus the probability is \(\frac{99}{100}\).
        \item staying always results in \(\frac{1}{100}.\) swapping gives probability 
        \begin{equation}
            \frac{99}{100} \frac{1}{99-k}
        \end{equation}
        where k is the number of doors opened. the probability is concentrated into \(99-k\) doors so the probability is distributed \(\frac{1}{99-k}\) ways.
        \item yes. the probability in the equation above is greater than \(\frac{1}{100} \) for all \(k\geq 1\)
    \end{enumerate}
    \item \begin{enumerate}
        \item since each person removes a month from the possibilities of the next persons birthday we see the probability is \begin{equation}
            \frac{12!}{(12-n)!(12)^n}
        \end{equation}
        for \(n\leq 12\). then 0 for \(n>12\). 
        \item the probability is \begin{equation}
            1-\frac{12!}{(12-n)!(12)^n}
        \end{equation}.
        each person removes a 1/12 probability that the next person has of colliding. 
        \item it takes 5 people. p(same birth month with 5 people) = .61806
        \item the probability 3 birthdays collide is \(\frac{1}{12}^2\). there should be approx \({n\choose 3}\) ways fo this to happen, thus
        \begin{equation}
            P_3(n) = \frac{{n\choose 3}}{144}
        \end{equation}
    \end{enumerate}

\end{enumerate}

\end{document}